\documentclass[10pt,a4paper]{amsart}
\usepackage{amsmath,amssymb,amsthm,graphicx}
\usepackage{verbatim}
\newenvironment{reference}{\comment}{\endcomment}
%\newenvironment{reference}{}{}
\newenvironment{slogan}{\comment}{\endcomment}
\newenvironment{history}{\comment}{\endcomment}

% For commutative diagrams we use Xy-pic
\usepackage[all]{xy}

% We use 2cell for 2-commutative diagrams.
\xyoption{2cell}
% Theorem environments.
%
\theoremstyle{plain}
\newtheorem{theorem}[subsection]{Theorem}
\newtheorem{proposition}[subsection]{Proposition}
\newtheorem{lemma}[subsection]{Lemma}

\theoremstyle{definition}
\newtheorem{definition}[subsection]{Definition}
\newtheorem{example}[subsection]{Example}
\newtheorem{exercise}[subsection]{Exercise}
\newtheorem{situation}[subsection]{Situation}

\theoremstyle{remark}
\newtheorem{remark}[subsection]{Remark}
\newtheorem{remarks}[subsection]{Remarks}

\numberwithin{equation}{subsection}

% Macros
%
\def\lim{\mathop{\mathrm{lim}}\nolimits}
\def\colim{\mathop{\mathrm{colim}}\nolimits}
\def\Spec{\mathop{\mathrm{Spec}}}
\def\Hom{\mathop{\mathrm{Hom}}\nolimits}
\def\Ext{\mathop{\mathrm{Ext}}\nolimits}
\def\SheafHom{\mathop{\mathcal{H}\!\mathit{om}}\nolimits}
\def\SheafExt{\mathop{\mathcal{E}\!\mathit{xt}}\nolimits}
\def\Sch{\mathit{Sch}}
\def\Mor{\mathop{\mathrm{Mor}}\nolimits}
\def\Ob{\mathop{\mathrm{Ob}}\nolimits}
\def\Sh{\mathop{\mathit{Sh}}\nolimits}
\def\NL{\mathop{N\!L}\nolimits}
\def\CH{\mathop{\mathrm{CH}}\nolimits}
\def\proetale{{pro\text{-}\acute{e}tale}}
\def\etale{{\acute{e}tale}}
\def\QCoh{\mathit{QCoh}}
\def\Ker{\mathop{\mathrm{Ker}}}
\def\Im{\mathop{\mathrm{Im}}}
\def\Coker{\mathop{\mathrm{Coker}}}
\def\Coim{\mathop{\mathrm{Coim}}}

% Boxtimes
%
\DeclareMathSymbol{\boxtimes}{\mathbin}{AMSa}{"02}

%
% Macros for moduli stacks/spaces
%
\def\QCohstack{\mathcal{QC}\!\mathit{oh}}
\def\Cohstack{\mathcal{C}\!\mathit{oh}}
\def\Spacesstack{\mathcal{S}\!\mathit{paces}}
\def\Quotfunctor{\mathrm{Quot}}
\def\Hilbfunctor{\mathrm{Hilb}}
\def\Curvesstack{\mathcal{C}\!\mathit{urves}}
\def\Polarizedstack{\mathcal{P}\!\mathit{olarized}}
\def\Complexesstack{\mathcal{C}\!\mathit{omplexes}}
% \Pic is the operator that assigns to X its picard group, usage \Pic(X)
% \Picardstack_{X/B} denotes the Picard stack of X over B
% \Picardfunctor_{X/B} denotes the Picard functor of X over B
\def\Pic{\mathop{\mathrm{Pic}}\nolimits}
\def\Picardstack{\mathcal{P}\!\mathit{ic}}
\def\Picardfunctor{\mathrm{Pic}}
\def\Deformationcategory{\mathcal{D}\!\mathit{ef}}
\usepackage{cleveref}
\crefname{originaldefinition}{original definition}{original definitions}
\Crefname{originaldefinition}{Original definition}{Original definitions}
\makeatletter
\newlabel{definition-ord}{{02MD}{source}{}{sourcepointer-1}{}}
\newlabel{definition-ord@cref}{{[originaldefinition][0][]02MD}{[1][0][]source}{}{}{}}
\makeatother
\crefname{originallemma}{original lemma}{original lemmas}
\Crefname{originallemma}{Original lemma}{Original lemmas}
\makeatletter
\newlabel{lemma-finite-extension-dim-1}{{02MJ}{source}{}{sourcepointer-2}{}}
\newlabel{lemma-finite-extension-dim-1@cref}{{[originallemma][0][]02MJ}{[1][0][]source}{}{}{}}
\makeatother
\crefname{originallemma}{original lemma}{original lemmas}
\Crefname{originallemma}{Original lemma}{Original lemmas}
\makeatletter
\newlabel{lemma-finite-in-codim-1}{{02MA}{source}{}{sourcepointer-3}{}}
\newlabel{lemma-finite-in-codim-1@cref}{{[originallemma][0][]02MA}{[1][0][]source}{}{}{}}
\makeatother
\crefname{originallemma}{original lemma}{original lemmas}
\Crefname{originallemma}{Original lemma}{Original lemmas}
\makeatletter
\newlabel{lemma-quasi-finite-open-integral-closure}{{00QB}{source}{}{sourcepointer-4}{}}
\newlabel{lemma-quasi-finite-open-integral-closure@cref}{{[originallemma][0][]00QB}{[1][0][]source}{}{}{}}
\makeatother
\title{Valuation bounds in a quasi-finite extension of a one-dimensional local domain, with equality exactly in the finite case}
\author{The Stacks Project Authors}
\date{stacks-project:algebra.tex:local-snapshot-20261012; GNU FDL 1.2 or later}
\begin{document}
\maketitle
\noindent\textit{Editorial note.} All delivered mathematical text is the author's own, copied from the pinned Stacks Project chapter file algebra.tex; no statement or proof step was rewritten, reordered or removed.\par\medskip
\noindent\textit{Editorial note.} Preamble provenance, as required for this source: the frozen snapshot stages the chapter's own preamble file beside it (preamble.tex, 217 lines, 8116 bytes, the same shared Stacks-wide preamble the published chow and etale-cohomology cards use); the first line of the pinned chapter file is \textbackslash{}input\{preamble\}, so no preamble text lives inside the chapter itself. The card takes that file's lines 8-18 (the verbatim package and the three \textbackslash{}comment-based environments reference/slogan/history at lines 9-12, section 13.2, together with the Xy-pic setup) and its lines 154 to the runtime end of the file (the native theorem declarations and mathematical macros). Not one character was changed, no package was added, no macro was defined, and the excerpt serves only as preamble, never as a span coordinate.\par\medskip
\noindent\textit{Editorial note.} The same-chapter pointers are located at their \textbackslash{}label lines inside the pinned chapter file; the cross-chapter pointers are located in the project's own tag list staged in the snapshot (sources/stacks-chow/tags.txt), at the verbatim lines carrying their canonical tags.\par\medskip
\noindent\textit{Editorial note.} The counted statement and proof refer to same-chapter statements and to the definition of the order function; the latter is recorded with kind definition in the navigation stubs, and the arguments consuming them stay inside the counted span.\par\medskip
\noindent\textit{Editorial note.} The counted spans contain no \textbackslash{}cite, so the card carries no bibliography.\par\medskip
\noindent\textit{Allowed standard prerequisites (editorial).}\begin{itemize}
\item Finiteness detected on the integral closure. Invoked in the counted proof (lines 31205-31209). Replaces the extension by a finite subalgebra of the integral closure, which is where the comparison of orders can be made.
\item The order function of a one-dimensional local domain. Used in the inequality and in the equality clause (lines 31193-31201). Supplies the lengths appearing on both sides of the displayed estimate.
\end{itemize}\medskip
\ifdefined\hypertarget\hypertarget{sourcepointer-1}{}\fi
\noindent\textit{Source reference.} Original definition 02MD: sources/stacks-algebra/pinned-original-source.tex, source lines 30020--30020. Navigation pointer for \textbackslash{}ref\{definition-ord\} used inside the counted span (line 31200); the argument that consumes it is retained in the span.\par\medskip
\ifdefined\hypertarget\hypertarget{sourcepointer-2}{}\fi
\noindent\textit{Source reference.} Original lemma 02MJ: sources/stacks-algebra/pinned-original-source.tex, source lines 30255--30255. Navigation pointer for \textbackslash{}ref\{lemma-finite-extension-dim-1\} used inside the counted span (line 31215); the argument that consumes it is retained in the span.\par\medskip
\ifdefined\hypertarget\hypertarget{sourcepointer-3}{}\fi
\noindent\textit{Source reference.} Original lemma 02MA: sources/stacks-algebra/pinned-original-source.tex, source lines 27820--27820. Navigation pointer for \textbackslash{}ref\{lemma-finite-in-codim-1\} used inside the counted span (line 31205); the argument that consumes it is retained in the span.\par\medskip
\ifdefined\hypertarget\hypertarget{sourcepointer-4}{}\fi
\noindent\textit{Source reference.} Original lemma 00QB: sources/stacks-algebra/pinned-original-source.tex, source lines 31123--31123. Navigation pointer for \textbackslash{}ref\{lemma-quasi-finite-open-integral-closure\} used inside the counted span (line 31207); the argument that consumes it is retained in the span.\par\medskip
\begin{lemma}
\label{lemma-quasi-finite-extension-dim-1}
Let $A \subset B$ be an extension of domains. Assume
\begin{enumerate}
\item $A$ is a local Noetherian ring of dimension $1$,
\item $A \to B$ is of finite type, and
\item the induced extension $L/K$ of fraction fields is finite.
\end{enumerate}
Then $B$ is semi-local.
Let $x \in \mathfrak m_A$, $x \not = 0$.
Let $\mathfrak m_i$, $i = 1, \ldots, n$
be the maximal ideals of $B$.
Then
$$
[L : K]\text{ord}_A(x)
\geq
\sum\nolimits_i
[\kappa(\mathfrak m_i) : \kappa(\mathfrak m_A)]
\text{ord}_{B_{\mathfrak m_i}}(x)
$$
where $\text{ord}$ is defined as in Definition \ref{definition-ord}.
We have equality if and only if $A \to B$ is finite.
\end{lemma}

\begin{proof}
The ring $B$ is semi-local by Lemma \ref{lemma-finite-in-codim-1}.
Let $B'$ be the integral closure of $A$ in $B$. By
Lemma \ref{lemma-quasi-finite-open-integral-closure}
we can find a finite $A$-subalgebra $C \subset B'$ such that
on setting $\mathfrak n_i = C \cap \mathfrak m_i$ we have
$C_{\mathfrak n_i} \cong B_{\mathfrak m_i}$ and the primes
$\mathfrak n_1, \ldots, \mathfrak n_n$ are pairwise distinct.
The ring $C$ is semi-local by Lemma \ref{lemma-finite-in-codim-1}.
Let $\mathfrak p_j$, $j = 1, \ldots, m$ be the other maximal
ideals of $C$ (the ``missing points''). By
Lemma \ref{lemma-finite-extension-dim-1} we have
$$
\text{ord}_A(x^{[L : K]}) =
\sum\nolimits_i
[\kappa(\mathfrak n_i) : \kappa(\mathfrak m_A)]
\text{ord}_{C_{\mathfrak n_i}}(x)
+
\sum\nolimits_j
[\kappa(\mathfrak p_j) : \kappa(\mathfrak m_A)]
\text{ord}_{C_{\mathfrak p_j}}(x)
$$
hence the inequality follows. In case of equality we conclude that
$m = 0$ (no ``missing points''). Hence $C \subset B$ is an inclusion
of semi-local rings inducing a bijection on maximal ideals and
an isomorphism on all localizations at maximal ideals. So if $b \in B$,
then $I = \{x \in C \mid xb \in C\}$ is an ideal of $C$ which is not
contained in any of the maximal ideals of $C$, and hence $I = C$,
hence $b \in C$. Thus $B = C$ and $B$ is finite over $A$.
\end{proof}

\end{document}
